% ---------------------------------------------------------------------------
% ACT C -- BATTERY MODULE UNDER A TAKEOFF PULSE -- one-page customer-facing
% result sheet, written under the DEMO STANDARD v2 output rules R1-R10.
%
% STATE OF THIS SHEET: the confidence section and the module results are both
% complete and every number in both is measured. The module was solved over
% the full 900 s at a 0.5 s step; the eight per-cell histories, the spread,
% the peak table and the energy closure are all read from the solution files
% on disk by a reader that carries its own planted control.
%
% ONE FRAME REMAINS RESERVED AND DELIBERATELY EMPTY: the step-size comparison.
% The module exists at a second, halved step of 0.25 s, built by the same
% builder and identical in every other respect, but it has NOT been solved.
% A step-independence claim needs both arms. Nothing is inferred, drawn or
% interpolated from one step size into that frame.
%
% Every user-visible number traces to a source artifact by file and line.
% That mapping lives in README_SOURCES.md in this directory; it is not
% repeated here, so that this file carries no internal identifiers at all.
%
% The two curves of Figure 1 and the one curve of Figure 2 are generated
% directly from the solved fields, not drawn by hand.
%
% Self-contained: no external figure file, no packages beyond a base TeX Live
% install (no tikz/pgfplots on this box). Vector output at any zoom.
% Compile: pdflatex ACT_C_battery_module_sheet.tex
% ---------------------------------------------------------------------------
\documentclass[9pt,a4paper]{article}
\usepackage[margin=10mm,top=5mm,bottom=3mm]{geometry}
\usepackage{booktabs}
\usepackage{array}
\usepackage{xcolor}
\usepackage{graphicx}
\usepackage[T1]{fontenc}

\definecolor{rulegrey}{rgb}{0.45,0.45,0.45}
\definecolor{boxgrey}{rgb}{0.96,0.96,0.94}
\definecolor{holdgrey}{rgb}{0.62,0.62,0.62}
\definecolor{holdfill}{rgb}{0.975,0.975,0.975}
\definecolor{keyfill}{rgb}{0.93,0.945,0.965}

\renewcommand{\footnotesize}{\fontsize{6.9}{7.5}\selectfont}
\renewcommand{\scriptsize}{\fontsize{6.4}{7.0}\selectfont}
\renewcommand{\tiny}{\fontsize{5.6}{6.1}\selectfont}
\renewcommand{\arraystretch}{0.96}
\setlength{\parindent}{0pt}
\setlength{\parskip}{0.4pt}
\pagestyle{empty}

\newcommand{\rolesig}[1]{\textbf{\footnotesize#1}}

% A reserved, visibly empty frame. Nothing is drawn inside it, deliberately.
\newcommand{\reserved}[3]{%
  \fcolorbox{holdgrey}{holdfill}{\begin{minipage}[t][#2][c]{#1}
    \centering{\tiny\color{holdgrey}#3}
  \end{minipage}}}

\begin{document}
\fontsize{7.8}{8.7}\selectfont
\enlargethispage{26pt}

% ---------------------------------------------------------------- title -----
{\fontsize{13.5}{15}\selectfont\bfseries Battery module under a takeoff power pulse}\\[1pt]
{\fontsize{8.6}{9.8}\selectfont Eight air-cooled cells, a 60-second high-power
takeoff phase, then cruise to 900 seconds}\\[2pt]
{\color{rulegrey}\rule{\textwidth}{0.9pt}}\\[1pt]
{\footnotesize\textbf{Solver:} OpenFOAM \texttt{chtMultiRegionFoam}, transient
conduction in the solid module, convective cooling on the channel faces.
The channels are not resolved as a fluid region. Solved over the full 900
seconds. The one comparison not yet available is marked as such.}

% -------------------------------------------------- 1. confidence beat ------
\vspace{1pt}
\textbf{1. Before spending anything on your module: what the transient
conduction machinery is worth}

\vspace{0.5pt}
{\footnotesize A three-dimensional transient conduction problem with a
closed-form answer. A uniform block, initially hot, cooled on its faces by a
fixed convective coefficient. That is the same physics as a cell losing heat to
a cooling channel. Three independent readings are compared against it.
Temperatures are the dimensionless excess
$\theta = (T-T_{\infty})/(T_{0}-T_{\infty})$ at $Bi = 1.0$, $Fo = 0.2$. Success
criteria and tolerances are fixed in advance.}

\vspace{1.5pt}
{\fontsize{6.7}{7.7}\selectfont
\begin{tabular}{l rr r r r r r}
\multicolumn{8}{l}{\textbf{Table 1 --- Solved against the exact conduction solution.} Three mesh levels each ($20$, $40$, $80$ per side; finest $512\,000$ cells).}\\
\toprule
Reading of the solution & Exact & Solved & Deviation & Tolerance & Grid-refine. & Observed & Refine. \\
(dimensionless $\theta$) & value & (finest) & exact, \% & used, \% & uncert., \% & order & ratio \\
\midrule
Average through the block        & 0.6175896496 & 0.6176037117 & $+0.00228$ & 45.5 & 0.0021 & 1.9996 & 3.999 \\
Temperature at the block centre  & 0.8591137894 & 0.8591831159 & $+0.00807$ & 32.3 & 0.0104 & 2.0006 & 4.002 \\
Temperature at a cooled face     & 0.5814449853 & 0.5815192620 & $+0.01277$ & 42.6 & 0.0154 & 2.0130 & 4.036 \\
\bottomrule
\end{tabular}}

\vspace{1.5pt}
{\footnotesize\textbf{Verified against the exact conduction solution to within
$0.013\,\%$} --- the largest of the three deviations, not an average and not a
tolerance. Each used under half its allowed tolerance (worst $45.5\,\%$), at a
measured order of accuracy of $2.00$ against a theoretical $2.00$:
\textbf{second-order accurate by measurement, not assertion.}}

% ---------------------------------------- 2. what the module is -------------
\vspace{1.5pt}
{\color{rulegrey}\rule{\textwidth}{0.5pt}}\\[1pt]
\textbf{2. The module, the duty cycle, and the eight cells}

\vspace{1.5pt}
\begin{minipage}[t]{0.35\textwidth}
{\fontsize{6.8}{7.8}\selectfont
\begin{tabular}{l r l}
\toprule
Quantity & Value & Unit \\
\midrule
Cells in the stack                  & 8      & -- \\
Cell, along flow $\times$ thick     & 100$\times$30 & mm \\
Cooling channel gap                 & 3      & mm \\
Stack height, overall               & 261    & mm \\
Depth modelled                      & 1.0    & m \\
\midrule
Takeoff heat, per cell, for 60 s    & 15     & W \\
Cruise heat, per cell, to 900 s     & 4      & W \\
Initial and coolant temperature     & 293    & K \\
\midrule
Heat in, takeoff phase              & 7\,200  & J \\
Heat in, cruise phase               & 26\,880 & J \\
Heat in, total                      & 34\,080 & J \\
\bottomrule
\end{tabular}}
\end{minipage}\hfill
\begin{minipage}[t]{0.63\textwidth}
\vspace{-1mm}
{\footnotesize\textbf{Figure 1 --- Every cell, through the whole duty cycle.}
Rise above the $293\,\mathrm{K}$ coolant. Eight curves are drawn and fall on
two: the interior six agree to six decimals, and so do the end pair.}

\vspace{1mm}
\setlength{\unitlength}{1mm}\begin{picture}(76,25)(-8,-5.5)
\put(0,0){\line(1,0){60.0}}\put(0,0){\line(0,1){17.0}}
\put(0.00,0){\line(0,-1){0.9}}\put(0.00,-4.0){\makebox(0,0)[b]{\tiny 0}}
\put(20.00,0){\line(0,-1){0.9}}\put(20.00,-4.0){\makebox(0,0)[b]{\tiny 300}}
\put(40.00,0){\line(0,-1){0.9}}\put(40.00,-4.0){\makebox(0,0)[b]{\tiny 600}}
\put(60.00,0){\line(0,-1){0.9}}\put(60.00,-4.0){\makebox(0,0)[b]{\tiny 900}}
\put(0,0.00){\line(-1,0){0.9}}\put(-1.3,0.00){\makebox(0,0)[r]{\tiny 0.0}}
\put(0,7.56){\line(-1,0){0.9}}\put(-1.3,7.56){\makebox(0,0)[r]{\tiny 0.2}}
\put(0,15.11){\line(-1,0){0.9}}\put(-1.3,15.11){\makebox(0,0)[r]{\tiny 0.4}}
\put(4.00,0.00){\line(0,1){0.8}}
\put(4.00,1.40){\line(0,1){0.8}}
\put(4.00,2.80){\line(0,1){0.8}}
\put(4.00,4.20){\line(0,1){0.8}}
\put(4.00,5.60){\line(0,1){0.8}}
\put(4.00,7.00){\line(0,1){0.8}}
\put(4.00,8.40){\line(0,1){0.8}}
\put(4.00,9.80){\line(0,1){0.8}}
\put(4.00,11.20){\line(0,1){0.8}}
\put(4.00,12.60){\line(0,1){0.8}}
\put(4.00,14.00){\line(0,1){0.8}}
\put(4.00,15.40){\line(0,1){0.8}}
\put(4.00,16.80){\line(0,1){0.8}}
\put(4.00,17.60){\makebox(0,0)[b]{\tiny takeoff ends}}
\linethickness{0.5pt}
\qbezier(0.00,0.00)(0.83,0.94)(1.67,1.87)
\qbezier(1.67,1.87)(2.50,2.79)(3.33,3.72)
\qbezier(3.33,3.72)(4.17,4.19)(5.00,4.67)
\qbezier(5.00,4.67)(5.83,4.89)(6.67,5.10)
\qbezier(6.67,5.10)(7.50,5.31)(8.33,5.52)
\qbezier(8.33,5.52)(9.17,5.73)(10.00,5.93)
\qbezier(10.00,5.93)(10.83,6.14)(11.67,6.34)
\qbezier(11.67,6.34)(12.50,6.54)(13.33,6.74)
\qbezier(13.33,6.74)(14.17,6.94)(15.00,7.14)
\qbezier(15.00,7.14)(15.83,7.33)(16.67,7.53)
\qbezier(16.67,7.53)(17.50,7.72)(18.33,7.91)
\qbezier(18.33,7.91)(19.17,8.10)(20.00,8.29)
\qbezier(20.00,8.29)(20.83,8.48)(21.67,8.67)
\qbezier(21.67,8.67)(22.50,8.85)(23.33,9.03)
\qbezier(23.33,9.03)(24.17,9.21)(25.00,9.39)
\qbezier(25.00,9.39)(25.83,9.57)(26.67,9.75)
\qbezier(26.67,9.75)(27.50,9.93)(28.33,10.10)
\qbezier(28.33,10.10)(29.17,10.28)(30.00,10.45)
\qbezier(30.00,10.45)(30.83,10.62)(31.67,10.79)
\qbezier(31.67,10.79)(32.50,10.96)(33.33,11.12)
\qbezier(33.33,11.12)(34.17,11.29)(35.00,11.46)
\qbezier(35.00,11.46)(35.83,11.62)(36.67,11.78)
\qbezier(36.67,11.78)(37.50,11.94)(38.33,12.10)
\qbezier(38.33,12.10)(39.17,12.26)(40.00,12.42)
\qbezier(40.00,12.42)(40.83,12.57)(41.67,12.73)
\qbezier(41.67,12.73)(42.50,12.88)(43.33,13.03)
\qbezier(43.33,13.03)(44.17,13.19)(45.00,13.34)
\qbezier(45.00,13.34)(45.83,13.49)(46.67,13.63)
\qbezier(46.67,13.63)(47.50,13.78)(48.33,13.93)
\qbezier(48.33,13.93)(49.17,14.07)(50.00,14.22)
\qbezier(50.00,14.22)(50.83,14.36)(51.67,14.50)
\qbezier(51.67,14.50)(52.50,14.64)(53.33,14.78)
\qbezier(53.33,14.78)(54.17,14.92)(55.00,15.05)
\qbezier(55.00,15.05)(55.83,15.19)(56.67,15.33)
\qbezier(56.67,15.33)(57.50,15.46)(58.33,15.59)
\qbezier(58.33,15.59)(59.17,15.73)(60.00,15.86)
\linethickness{0.22pt}
\qbezier(0.00,0.00)(0.83,0.93)(1.67,1.86)
\qbezier(1.67,1.86)(2.50,2.75)(3.33,3.65)
\qbezier(3.33,3.65)(4.17,4.10)(5.00,4.54)
\qbezier(5.00,4.54)(5.83,4.71)(6.67,4.89)
\qbezier(6.67,4.89)(7.50,5.06)(8.33,5.23)
\qbezier(8.33,5.23)(9.17,5.39)(10.00,5.55)
\qbezier(10.00,5.55)(10.83,5.71)(11.67,5.87)
\qbezier(11.67,5.87)(12.50,6.02)(13.33,6.17)
\qbezier(13.33,6.17)(14.17,6.32)(15.00,6.47)
\qbezier(15.00,6.47)(15.83,6.61)(16.67,6.76)
\qbezier(16.67,6.76)(17.50,6.89)(18.33,7.03)
\qbezier(18.33,7.03)(19.17,7.17)(20.00,7.30)
\qbezier(20.00,7.30)(20.83,7.43)(21.67,7.56)
\qbezier(21.67,7.56)(22.50,7.69)(23.33,7.81)
\qbezier(23.33,7.81)(24.17,7.93)(25.00,8.05)
\qbezier(25.00,8.05)(25.83,8.17)(26.67,8.29)
\qbezier(26.67,8.29)(27.50,8.40)(28.33,8.52)
\qbezier(28.33,8.52)(29.17,8.63)(30.00,8.74)
\qbezier(30.00,8.74)(30.83,8.84)(31.67,8.95)
\qbezier(31.67,8.95)(32.50,9.05)(33.33,9.16)
\qbezier(33.33,9.16)(34.17,9.26)(35.00,9.36)
\qbezier(35.00,9.36)(35.83,9.45)(36.67,9.55)
\qbezier(36.67,9.55)(37.50,9.64)(38.33,9.74)
\qbezier(38.33,9.74)(39.17,9.83)(40.00,9.92)
\qbezier(40.00,9.92)(40.83,10.00)(41.67,10.09)
\qbezier(41.67,10.09)(42.50,10.18)(43.33,10.26)
\qbezier(43.33,10.26)(44.17,10.34)(45.00,10.43)
\qbezier(45.00,10.43)(45.83,10.50)(46.67,10.58)
\qbezier(46.67,10.58)(47.50,10.66)(48.33,10.74)
\qbezier(48.33,10.74)(49.17,10.81)(50.00,10.89)
\qbezier(50.00,10.89)(50.83,10.96)(51.67,11.03)
\qbezier(51.67,11.03)(52.50,11.10)(53.33,11.17)
\qbezier(53.33,11.17)(54.17,11.24)(55.00,11.30)
\qbezier(55.00,11.30)(55.83,11.37)(56.67,11.43)
\qbezier(56.67,11.43)(57.50,11.50)(58.33,11.56)
\qbezier(58.33,11.56)(59.17,11.62)(60.00,11.68)
\put(60.70,15.86){\makebox(0,0)[l]{\tiny cells 1, 8}}
\put(60.70,11.68){\makebox(0,0)[l]{\tiny cells 2--7}}
\put(30.00,-5.9){\makebox(0,0)[b]{\tiny time, s}}
\put(-7.0,8.50){\rotatebox{90}{\makebox(0,0)[b]{\tiny rise above coolant, K}}}
\end{picture}
\end{minipage}

% ---------------------------------------- 3. the headline -------------------
\vspace{1pt}
\colorbox{keyfill}{\begin{minipage}{0.988\textwidth}
\vspace{1pt}
{\footnotesize\textbf{Three things this result should change about how you size
the pack.} \textbf{(1) The takeoff pulse is not what sets your peak.} Only
$0.117\,\mathrm{K}$ of the rise has arrived when the takeoff phase ends.
Takeoff puts in $7\,200\,\mathrm{J}$; cruise puts in $26\,880\,\mathrm{J}$,
because it lasts fourteen times longer --- the cruise soak is the design driver.
\textbf{(2) The module is not uniform.} The two end cells are cooled on one
face, the interior cells on two. That gives the end cells twice the time
constant and twice the settled temperature. \textbf{(3) The window closes before the
cells that matter stop rising.} The end cells reach $0.420\,\mathrm{K}$ at
$900\,\mathrm{s}$, and are heading for roughly $0.74\,\mathrm{K}$. At that
point the channel removes only $22.1\,\mathrm{W}$ of the $32\,\mathrm{W}$
going in. \textbf{$0.420\,\mathrm{K}$ is where the hottest cells are when the
clock stops, not where they are going.}\\[1.5pt]
\textbf{Energy audit.} Storage and surface removal are measured separately.}

\vspace{0.5pt}
{\fontsize{6.7}{7.4}\selectfont
\begin{tabular}{l r r r r}
\toprule
Energy over $0$--$900\,\mathrm{s}$ & Heat in, J & Stored, J &
Removed, J & Closure \\
\midrule
Module, eight cells & 34\,080 & 20\,212 & 13\,818 & $\mathbf{99.85\,\%}$ \\
\bottomrule
\end{tabular}}
\vspace{1pt}
\end{minipage}}

\vspace{1.5pt}
\begin{minipage}[t]{0.325\textwidth}
{\fontsize{6.8}{7.9}\selectfont
\begin{tabular}{c r r r l}
\multicolumn{5}{l}{\textbf{Table 2 --- Peak per cell}}\\
\toprule
Cell & Peak & Rise & Time to & Uncer- \\
     & $^{\circ}$C & K & peak, s & tainty \\
\midrule
1 & 20.27 & 0.420 & 900 & n/q \\
2 & 20.16 & 0.309 & 900 & n/q \\
3 & 20.16 & 0.309 & 900 & n/q \\
4 & 20.16 & 0.309 & 900 & n/q \\
5 & 20.16 & 0.309 & 900 & n/q \\
6 & 20.16 & 0.309 & 900 & n/q \\
7 & 20.16 & 0.309 & 900 & n/q \\
8 & 20.27 & 0.420 & 900 & n/q \\
\bottomrule
\end{tabular}}

{\tiny Every cell peaks at $900\,\mathrm{s}$ because that is where the run
ends, not because it turns over. Hottest point anywhere:
$20.30\,^{\circ}\mathrm{C}$, at the insulated casing. ``n/q'' = not quantified
--- one mesh, one time step.}
\end{minipage}\hfill
\begin{minipage}[t]{0.325\textwidth}
{\footnotesize\textbf{Figure 2 --- Spread across the module}}

\vspace{0.5mm}
\setlength{\unitlength}{1mm}\begin{picture}(62,17)(-8,-5.5)
\put(0,0){\line(1,0){52.0}}\put(0,0){\line(0,1){10.0}}
\put(0.00,0){\line(0,-1){0.9}}\put(0.00,-4.0){\makebox(0,0)[b]{\tiny 0}}
\put(26.00,0){\line(0,-1){0.9}}
\put(52.00,0){\line(0,-1){0.9}}\put(52.00,-4.0){\makebox(0,0)[b]{\tiny 900}}
\put(0,0.00){\line(-1,0){0.9}}\put(-1.3,0.00){\makebox(0,0)[r]{\tiny 0.00}}
\put(0,4.17){\line(-1,0){0.9}}\put(-1.3,4.17){\makebox(0,0)[r]{\tiny 0.05}}
\put(0,8.33){\line(-1,0){0.9}}\put(-1.3,8.33){\makebox(0,0)[r]{\tiny 0.10}}
\linethickness{0.5pt}
\qbezier(0.00,0.00)(0.72,0.02)(1.44,0.04)
\qbezier(1.44,0.04)(2.17,0.09)(2.89,0.14)
\qbezier(2.89,0.14)(3.61,0.22)(4.33,0.29)
\qbezier(4.33,0.29)(5.06,0.38)(5.78,0.47)
\qbezier(5.78,0.47)(6.50,0.56)(7.22,0.65)
\qbezier(7.22,0.65)(7.94,0.74)(8.67,0.84)
\qbezier(8.67,0.84)(9.39,0.94)(10.11,1.04)
\qbezier(10.11,1.04)(10.83,1.15)(11.56,1.26)
\qbezier(11.56,1.26)(12.28,1.37)(13.00,1.48)
\qbezier(13.00,1.48)(13.72,1.59)(14.44,1.71)
\qbezier(14.44,1.71)(15.17,1.83)(15.89,1.94)
\qbezier(15.89,1.94)(16.61,2.07)(17.33,2.19)
\qbezier(17.33,2.19)(18.06,2.31)(18.78,2.44)
\qbezier(18.78,2.44)(19.50,2.57)(20.22,2.70)
\qbezier(20.22,2.70)(20.94,2.83)(21.67,2.96)
\qbezier(21.67,2.96)(22.39,3.09)(23.11,3.23)
\qbezier(23.11,3.23)(23.83,3.36)(24.56,3.50)
\qbezier(24.56,3.50)(25.28,3.64)(26.00,3.78)
\qbezier(26.00,3.78)(26.72,3.92)(27.44,4.06)
\qbezier(27.44,4.06)(28.17,4.20)(28.89,4.34)
\qbezier(28.89,4.34)(29.61,4.49)(30.33,4.63)
\qbezier(30.33,4.63)(31.06,4.78)(31.78,4.92)
\qbezier(31.78,4.92)(32.50,5.07)(33.22,5.22)
\qbezier(33.22,5.22)(33.94,5.37)(34.67,5.52)
\qbezier(34.67,5.52)(35.39,5.67)(36.11,5.82)
\qbezier(36.11,5.82)(36.83,5.97)(37.56,6.12)
\qbezier(37.56,6.12)(38.28,6.27)(39.00,6.42)
\qbezier(39.00,6.42)(39.72,6.57)(40.44,6.73)
\qbezier(40.44,6.73)(41.17,6.88)(41.89,7.04)
\qbezier(41.89,7.04)(42.61,7.19)(43.33,7.34)
\qbezier(43.33,7.34)(44.06,7.50)(44.78,7.65)
\qbezier(44.78,7.65)(45.50,7.81)(46.22,7.96)
\qbezier(46.22,7.96)(46.94,8.12)(47.67,8.27)
\qbezier(47.67,8.27)(48.39,8.43)(49.11,8.59)
\qbezier(49.11,8.59)(49.83,8.74)(50.56,8.90)
\qbezier(50.56,8.90)(51.28,9.05)(52.00,9.21)
\put(26.00,-5.9){\makebox(0,0)[b]{\tiny time, s}}
\put(-7.0,5.00){\rotatebox{90}{\makebox(0,0)[b]{\tiny spread, K}}}
\end{picture}

\vspace{-1mm}
{\tiny Hottest minus coldest cell. It grows monotonically to
$0.111\,\mathrm{K}$ and never turns over.}
\end{minipage}\hfill
\begin{minipage}[t]{0.315\textwidth}
{\fontsize{6.8}{7.9}\selectfont
\begin{tabular}{l r r}
\multicolumn{3}{l}{\textbf{Table 3 --- How far from settled}}\\
\toprule
 & End cells & Interior \\
 & 1, 8 & 2--7 \\
\midrule
Cooled area, m$^2$      & 0.100  & 0.200 \\
Time constant, s        & 1\,392 & 696 \\
Heading for, K$^{*}$    & 0.742  & 0.371 \\
Reached at 900 s, K     & 0.420  & 0.309 \\
Fraction of that        & 57\,\% & 83\,\% \\
Still rising at, K/h    & 1.00   & 0.46 \\
\bottomrule
\end{tabular}}

\vspace{0.5mm}
{\tiny $^{*}$A \emph{conservative upper} estimate: it assumes each cell is
thermally on its own, which is how this model is built. Busbars, casing
conduction or coolant cross-talk in a real pack would pull it down.}
\end{minipage}

% ------------------------------------------------------------ caveats -------
\vspace{1.5pt}
{\color{rulegrey}\rule{\textwidth}{0.5pt}}\\[1pt]
\begin{minipage}[t]{0.615\textwidth}
\vspace{0pt}
\colorbox{boxgrey}{\begin{minipage}{0.975\textwidth}
{\footnotesize\textbf{Read this before you use anything on this sheet}\\[1pt]
$\bullet$ \textbf{The cooling coefficient is estimated, not measured, and your
peak scales inversely with it.} It is calculated from the stated channel flow
($8\,\mathrm{m\,s^{-1}}$ in a $3\,\mathrm{mm}$ gap) and validated against no
test. Settled rise is heat per cell over coefficient times cooled area ---
\emph{exactly} inverse, so \textbf{a channel that cools half as well gives twice
the rise.}\\
$\bullet$ \textbf{The run ends before the module settles.} The end cells are at
$57\,\%$ of where they are heading, still climbing at $1\,\mathrm{K}$ per hour,
and approach roughly $0.74\,\mathrm{K}$ --- near double this sheet's peak. That
$0.74\,\mathrm{K}$ is a \emph{conservative upper} figure. The model gives each
cell its own cooling and nothing else. A real pack with busbars, casing
conduction or coolant cross-talk comes in under it.\\
$\bullet$ \textbf{Those two multiply, they do not sit side by side.} An
optimistic coefficient scales the whole curve up and the hottest cells have not
finished climbing it. \textbf{Do not design against $0.42\,\mathrm{K}$ as if it
were robust.}\\
$\bullet$ \textbf{The channels are not resolved as air flow} but as a convective
condition on the channel-facing surfaces. So Figure 2's spread is end-cell
geometry and nothing else. The six interior cells agree to six decimals, and so
do the two end cells. It excludes coolant heating along the channel, which this
model cannot see and which would separate all eight.\\
$\bullet$ Properties are representative, not measured for your cells.
Conductivity is equal in all directions. The stack ends are insulated with no
plenum, and there is no cell-to-cell contact path. \textbf{One mesh and one time
step, so no numerical error bar is quoted on any module number.} Section 1 says
nothing about the channel airflow.}
\end{minipage}}
\end{minipage}\hfill
\begin{minipage}[t]{0.37\textwidth}
\vspace{0pt}
{\footnotesize\textbf{Compute}\\[1pt]
{\fontsize{6.8}{7.9}\selectfont
\begin{tabular}{l r r}
\toprule
Stage & Time & Cost \\
\midrule
Confidence check, section 1 & 2\,h\,05\,min & $\$0.11$ \\
The module run              & $1\,\mathrm{s}$ & under $\$0.01$ \\
\bottomrule
\end{tabular}}\\[2pt]
The whole $900\,\mathrm{s}$ of module physics costs under a penny. Costs are
derived at $\$0.0513$ per core-hour, a rate stated by the machine's owner
rather than metered here.}

\vspace{2pt}
% ---- THE ONE FRAME STILL RESERVED ----------------------------------------
% FILLED BY: the same module re-solved at half the time step, 0.25 s; peak
% temperature and time-to-peak compared between the two, difference in K and s.
% Both cases exist and are built, differing in that one entry and nothing else.
% The 0.25 s case has NOT been solved. Nothing may be written here until it is:
% a step-independence claim needs both arms, and a comparison drawn from one
% step size would be worse than an empty frame.
\reserved{0.97\textwidth}{11.5mm}{\textbf{Time-step check}\\[1.5pt]
Peak temperature and time to peak at time steps of $0.5\,\mathrm{s}$ and
$0.25\,\mathrm{s}$. The second is not yet available; nothing is drawn here.}
\end{minipage}

% ------------------------------------------------------------- voices -------
\vspace{1pt}
{\color{rulegrey}\rule{\textwidth}{0.5pt}}

\vspace{0.5pt}
{\footnotesize
\rolesig{Lead Researcher.} The exact-solution check comes first on purpose, and
it licenses everything after it. The module result is clear-cut: the event you
named the case after is not the event that sizes it.

\rolesig{Lead Engineer.} Size this pack against the cruise soak, not the takeoff
pulse; the pulse is a fifth of the energy. The two end cells are your problem,
and a channel outboard of cells 1 and 8 removes that asymmetry outright.

\rolesig{Lead Numericist.} Section 1 is a genuine verification. The module run is
not, and I will not dress it as one: one mesh, one time step, and no numerical
error bar on any module number.\vspace{-2pt}}

\end{document}
